\section{Sequential interventions and the longitudinal g-formula: proofs}
\label{sec:annex-p16-sequential-g-formula}

\begin{annexitem}{annex:p16.sequential-policy-semantics}{thm:p16-sequential-policy-semantics}{Finite table evaluation and attained histories}
\begin{proof}
There are finitely many formal histories and each selector has a finite
nonempty value set. Give the enlarged outcome \((\omega,r)\) mass
\(q(\omega)\prod_{k,h}\pi_k(r_{k,h}\mid h)\). Repeated finite summation
shows that this is a mass function, with the specified independent marginals.
The baseline is already fixed. Once the history through \(k\) has been
evaluated, its index selects one existing entry of the finite selector
table, thereby determining \(A_k^\pi\). That action extends the index of
the existing potential covariate table for the next stage. After the last
action the potential response table supplies \(Y^\pi\). Induction gives
existence and uniqueness in exactly \(K+1\) treatment stages.

If factual and policy actions agree up to \(k-1\), consistency gives
\(L_j[\bar A_{j-1}]=L_j\) at every \(1\leq j\leq k\), and their
baselines are identical, so the two histories
agree. If the selector at this factual history also chooses the factual
\(A_k\), agreement extends by another action. Starting with their common
baseline proves the partial-history assertion, and terminal consistency
gives \(Y[\bar A_K]=Y\). This is a pointwise assertion even at an original
or enlarged outcome of mass zero.

Explicitly define the reduced deterministic rule by
\[
 g'_0(l_0)=g_0(l_0,\varnothing),\qquad
 g'_k(\bar l_k)=g_k\bigl(\bar l_k,
     (g'_0(l_0),\ldots,g'_{k-1}(\bar l_{k-1}))\bigr).
\]
These are total functions on finite covariate-history sets. At stage zero
the prescribed actions agree. If they agree through the preceding stage,
the definition makes their next actions agree whenever their covariate
histories agree. The potential covariate tables therefore return the same
next values. Induction proves equality of all trajectories and responses.
For a factual history, the same induction proves the equivalence of its
two adherence conditions. When the original rules have no covariate
arguments, their recursively chosen values are constants from time zero,
so they define a single fixed sequence. Finally the number of fixed
sequences is the size of a Cartesian product of the action sets, proving
the count used in the body.
\end{proof}
\end{annexitem}

\begin{annexitem}{annex:p16.sequential-exchangeability-bridge}{thm:p16-sequential-exchangeability-bridge}{Joint future events, rather than separate marginals}
\begin{proof}
Fix \(g,k\) and a positive compatible observed history \(h_k\).
Specify a complete possible value
\((l_{k+1},\ldots,l_K,y)\) of its future policy covariates and response.
Starting at \(h_k\), recursively apply the rules to these proposed future
covariates to obtain a fixed continuation \((a_k,\ldots,a_K)\); attach
the actions already recorded in \(h_k\). This gives one fixed sequence
\(\bar a\), independently of the individual's original outcome.

On the conditioning event \(H_k=h_k\), the policy's earlier actions and
covariates agree with the recorded ones by consistency and compatibility.
Induction over the proposed continuation therefore gives the event equality
\[
 \{(L_{k+1}^g,\ldots,L_K^g,Y^g)=(l_{k+1},\ldots,l_K,y)\}
 =\{T_k(\bar a)=(l_{k+1},\ldots,l_K,y)\}
 \quad\hbox{within }\{H_k=h_k\}.
\]
For the forward implication the attained policy actions are the computed
ones; for the reverse implication each next table entry has the proposed
value, forcing the same next action. This also proves the equality at the
response. It holds irrespective of the factual value of the current
\(A_k\). By (SE--J), the indicator of the right-hand event is independent
of that factual action conditional on \(H_k=h_k\). Thus every cell of the
future policy tuple has the required conditional factorization. These cells
partition the finite tuple space, so this proves its joint independence.
Taking the marginal on the response gives (SE--Y).

Conversely a constant sequence is a deterministic policy. Its future tuple
is exactly \(T_k(\bar a)\), and its compatible histories are those with
the fixed earlier actions. The asserted dynamic family therefore includes
every instance of (SE--J). The third counterexample in
\cref{thm:p16-sequential-failures}, proved below without using this result,
has the two separate marginal independences and fails the joint one.
Consequently the joint premise in the forward direction is substantive.
\end{proof}
\end{annexitem}

\begin{annexitem}{annex:p16.sequential-structural-bridge}{thm:p16-sequential-structural-bridge}{Private assignment randomness conditional on recorded history}
\begin{proof}
The stipulated order is topological: each equation uses earlier observed
coordinates and a background input. The existing existence and uniqueness
theorem applies. Fix all treatment coordinates to a sequence \(\bar a\).
Then every treatment output is a constant and no treatment disturbance
is used in any retained equation. Induction in this order shows that every
potential covariate and response in that world is a function of \(U\)
alone. The entire finite tuple \(\mathcal T\), over all sequences, is
therefore a function \(t(U)\).

Later treatment replacements cannot affect an earlier covariate, as its
equation and all earlier inputs are unchanged. Hence its world index may
be shortened to the earlier treatment prefix. If the factual treatment
prefix equals that fixed prefix, the factual earlier trajectory satisfies
all the retained and replaced equations up to that covariate. Uniqueness,
or direct induction, proves the sequential consistency identities. Terminal
consistency is also \cref{thm:p16-counterfactual-laws}.

In the factual world \(H_k\) is a function of \(U,E_0,\ldots,E_{k-1}\).
The stipulated independence makes \(E_k\) independent of the joint pair
\((\mathcal T,H_k)\). For \(p(H_k=h)>0\), finite summation yields
\[
\begin{aligned}
 &\mathbb P(\mathcal T=t,A_k=a\mid H_k=h)\\
 &\quad=\sum_e\mathbf1_{\{f_k(h,e)=a\}}
       \mathbb P(\mathcal T=t,E_k=e\mid H_k=h)\\
 &\quad=\mathbb P(\mathcal T=t\mid H_k=h)
       \sum_e\mathbf1_{\{f_k(h,e)=a\}}\mathbb P(E_k=e).
\end{aligned}
\]
Summing this equality over \(t\) identifies its final factor as
\(\mathbb P(A_k=a\mid H_k=h)\). This proves the claimed full conditional
independence and all its one-world marginals. Dependence between components
of \(U\) is irrelevant to this argument.

For a policy, replace each treatment equation by the selector indexed by
its earlier history and retain all other equations and \(U\). All new
parents precede the treatment in the stated order, so the new system is
acyclic. At a fixed enlarged outcome, compare this solution with the static
potential solution indexed by its attained complete action sequence. The
non-treatment equations and inputs agree, so induction identifies its
covariates and response with the corresponding potential-table entries.
The treatment equations are exactly the policy recursion. Its solution is
therefore \(O^\pi\). This construction proves the bridge under the given
mechanism restrictions; it does not infer those restrictions from \(p\).
\end{proof}
\end{annexitem}

\begin{annexitem}{annex:p16.sequential-support}{thm:p16-sequential-support}{Every requested denominator and every zero branch}
\begin{proof}
At stage zero, positive \(m_0(l_0)\) is exactly positive observed baseline
mass. Suppose a reached \(h\) has positive observed mass. By (RP), an action
of positive policy mass has \(e_k(a\mid h)>0\), whence
\[
 p(H_k=h,A_k=a)=p(H_k=h)e_k(a\mid h)>0.
\]
If a next covariate has positive transition mass then
\[
 p(H_{k+1}=(h,a,l))=p(H_k=h,A_k=a)r_k(l\mid h,a)>0.
\]
This proves the positivity invariant by induction and licenses every
observed transition used. The same last equality with response \(y\)
shows that a positive terminal branch has \(p(o)>0\).

Summing all propagated children of a reached history gives
\[
 m_k(h)\sum_a\pi_k(a\mid h)\sum_l r_k(l\mid h,a)=m_k(h).
\]
Only actions of positive policy mass enter, and each of their transition
kernels is normalized. Summing over the finitely many histories proves
normalization at the next stage. At the terminal stage sum over \(a,y\)
in the same way. This proves that (GF) is a mass function. If a baseline,
policy or transition factor first becomes zero, its branch has no mass
to distribute; no later kernel can contribute to that branch's mass.
Any normalized assignments at null arguments are therefore immaterial.

For equivalence with the observed-law formulation, first assume (RP).
Take a positive observed \(h_k\) with all earlier policy factors positive.
All its observed prefixes and transitions are positive, since the event
\(H_k=h_k\) is contained in each relevant prefix event. Following those
prefixes in the recursion gives positive \(m_k(h_k)\), so (RP) imposes
the displayed action requirement there. Conversely, assume the observed-law
formulation and follow the recursion from a positive baseline. At a reached
positive observed history its preceding policy factors are positive, so
the formulation supplies positivity of every selected action. The induction
in the first paragraph supplies positivity of each next reached history.
Thus the recursion never requests a missing denominator and (RP) holds.
All assertions depend on the observed law and the specified policy alone.
\end{proof}
\end{annexitem}

\begin{annexitem}{annex:p16.sequential-joint-g-formula}{thm:p16-sequential-joint-g-formula}{A path calculation using joint counterfactual exchangeability}
\begin{proof}
Fix a proposed complete trajectory \(o=(\bar l,\bar a,y)\). Let
\[
 E_o=\{L_0=l_0, L_j[\bar a_{j-1}]=l_j\ (1\leq j\leq K),
                         \ Y[\bar a]=y\}.
\]
For a fixed original outcome \(\omega\), policy evaluation gives this
trajectory if and only if \(\omega\in E_o\) and
\(R_{k,h_k}=a_k\) at all its proposed histories. The equivalence follows
by induction along the path. The entries used have distinct time indices,
so they are independent even when some history value sets happen to
coincide. Independence from \(\omega\) gives
\[
 \mathbb P(O^\pi=o)=\mathbb P(E_o)\prod_{k=0}^K\pi_k(a_k\mid h_k).
 \tag{1}
\]
If any policy factor is zero, both sides of (GF) assign this trajectory
zero without further conditional evaluation. Likewise a zero baseline
mass forces \(\mathbb P(E_o)=0\).

Consider a path with all policy factors and baseline mass positive. Put
\(r^{\mathrm{tail}}_k=(l_{k+1},\ldots,l_K,y)\), and, on a positive
observed \(h_k\), abbreviate
\[
 F_k=\mathbb P(T_k(\bar a)=r^{\mathrm{tail}}_k\mid H_k=h_k).
\]
Initially \(\mathbb P(E_o)=b(l_0)F_0\). At a reached history (RP)
licenses conditioning additionally on \(A_k=a_k\). Joint exchangeability
therefore gives
\[
 F_k=\mathbb P(T_k(\bar a)=r^{\mathrm{tail}}_k
                                      \mid H_k=h_k,A_k=a_k).
 \tag{2}
\]
For \(k<K\), consistency identifies the first coordinate
\(L_{k+1}[\bar a_k]\) with the factual \(L_{k+1}\) on that event.
If \(r_k(l_{k+1}\mid h_k,a_k)=0\), equation (2) is zero, hence
\(\mathbb P(E_o)=0\); no remaining division is needed. Otherwise the
next history is positive by the support lemma. Partitioning (2) by its
first coordinate and applying the definition of finite conditioning gives
\[
 F_k=r_k(l_{k+1}\mid h_k,a_k)F_{k+1}.               \tag{3}
\]
At the last stage, (2) and terminal consistency give
\(F_K=r_K(y\mid h_K,a_K)\). Repeated substitution proves
\[
 \mathbb P(E_o)=b(l_0)\prod_{k<K}r_k(l_{k+1}\mid h_k,a_k)
                                      r_K(y\mid h_K,a_K)
\]
on every contributing path, including a path terminated at its first zero
transition. Substitution into (1) proves the full trajectory formula.

Summing it over future variables proves that \(m_k\) is the actual policy
history law. This justifies the reachable-history interpretation of (RP).
Summing over actions and covariates gives the response law; multiplication
by the finite response values and summation gives its mean. A deterministic
policy assigns zero to every nonprescribed action, so its action sums
collapse to the recursively specified choices. For static policies these
are the same constants at every covariate path, giving the stated formula.
Every surviving term is determined by \(p\) and \(\pi\), with licensed
denominators. Two potential-table models in the stipulated class sharing
these inputs have identical targets. No independence between their entire
collections of counterfactual worlds was used.

The operational support condition implies the stated actual-history
condition because the proved marginal identity is
\(\mathbb P(H_k^\pi=h)=m_k(h)\). For the converse, assume the
actual-history condition and suppose the operational recursion has a
first failed support request at time \(k\) and history \(h\).
The baseline law agrees in the two constructions. More generally the
path proof (1)--(3), stopped upon reaching \(L_k\), proves equality
of their history masses through that stage using only successful requests
at times \(0,\ldots,k-1\): replace \(E_o\) by the event prescribing
just the baseline and potential covariates through \(L_k\), project
(SE--J) to the required remaining covariate coordinates, and finish with
the last covariate transition instead of a response transition. Empty
remaining tuples have conditional probability one. No request at or
after \(k\) enters this truncated argument. Thus the positive reached
mass at the proposed first failure is the positive actual policy mass.
The selected action has positive policy probability, so the assumed
actual-history condition supplies both required positive observed masses,
a contradiction. Finiteness ensures that any operational failure has
a first time. This proves the claimed equivalence.
\end{proof}
\end{annexitem}

\begin{annexitem}{annex:p16.sequential-outcome-g-formula}{thm:p16-sequential-outcome-g-formula}{Backward conditioning for one policy's response}
\begin{proof}
Fix \(g\) and \(\varphi\). At a positive reached history, which is
necessarily compatible with \(g\), define
\(F_k(h)=\mathbb E[\varphi(Y^g)\mid H_k=h]\).
Since (SE--Y) is independence of the full finite response, it also holds
after applying \(\varphi\). The selected action has positive probability,
so dividing the corresponding conditional factorization gives
\[
 F_k(h)=\mathbb E[\varphi(Y^g)\mid H_k=h,A_k=g_k(h)].       \tag{4}
\]
At \(k=K\), the conditioning event describes a fully adhering factual
trajectory. Policy consistency gives \(Y^g=Y\) there, so
\[
 F_K(h)=\sum_y\varphi(y)r_K(y\mid h,g_K(h)).
\]
For \(k<K\), partition the positive event in (4) by \(L_{k+1}\).
Terms with zero transition probability have zero contribution and require
no future conditional mean. Every other next history is reached and
positive. Finite conditional total expectation gives
\[
 F_k(h)=\sum_{l:r_k(l\mid h,g_k(h))>0}
                  r_k(l\mid h,g_k(h))F_{k+1}(h,g_k(h),l).   \tag{5}
\]
This is a backward recurrence using only the observed transitions and
terminal conditional outcome mean. At baseline, total expectation gives
\(\mathbb E[\varphi(Y^g)]=\sum_{l_0}b(l_0)F_0(l_0)\).
Expanding (5) finitely to its terminal expression yields exactly
\(\sum_o\varphi(y)q_g(o)\). Taking indicators of response values proves
the response distribution claim. This calculation never asserted that a
future potential covariate is exchangeable with an earlier action. Its
observed intermediate factors need not be the corresponding causal
covariate transitions. The joint theorem's stronger claim is therefore
not a hidden premise of this proof.

For the stated limitation, take independent fair bits \(U,E\), constant
\(L_0\), and equations \(A_0=U,A_1=E,Y=0\). In one model put
\(L_1=A_0\oplus U\); in another put \(L_1=0\). Both have the same
observed law: \(L_1=Y=0\) and independent fair treatments. Every policy
response is zero pointwise, so (SE--Y) holds for every deterministic policy.
For the static policy \((0,0)\), the recursive support test succeeds:
both required action probabilities on its reached observed history are
\(1/2\). Its g-formula outcome is zero in both models, as required.
But \(L_1[0]=U\) in the first model and zero in the second. Their true
intervened covariate laws therefore differ. In the first model the actual
history \(A_0^g=0,L_1^g=1\) has mass \(1/2\) but no observed mass;
(SE--J) fails there at the earlier assignment mechanism. Thus the
outcome-only argument neither identifies the covariate law nor licenses
using the joint theorem's actual-history interpretation without its joint
exchangeability premise.
\end{proof}
\end{annexitem}

\begin{annexitem}{annex:p16.sequential-ice}{thm:p16-sequential-ice}{Backward summation of the policy probability tree}
\begin{proof}
At a reached \(h_K\), summing the terminal response contribution
\(\varphi(y)\pi_K(a\mid h_K)r_K(y\mid h_K,a)\) over \(a,y\)
gives exactly \(V_K(h_K)\). Suppose \(V_{k+1}(h')\) equals the sum of
all continuation contributions from each reached next history. From \(h_k\)
each next edge has mass multiplier
\(\pi_k(a\mid h_k)r_k(l\mid h_k,a)\). Multiplying the inductively
computed continuation by this multiplier and summing over \(a,l\)
is precisely the stated definition of \(V_k(h_k)\). This induction
computes the complete future path sum. Summing \(b(l_0)V_0(l_0)\)
over baseline states therefore gives the (GF) sum. Finiteness permits
every change of summation order. There are only \(K+1\) stages and
finitely many histories per stage, so the calculation terminates.

On a selected positive \((h,a)\), the observed conditional law of
\(H_{k+1}\) is obtained by adjoining \(a\) and drawing \(L_{k+1}\)
with mass \(r_k(\cdot\mid h,a)\). Its conditional expectation of
\(V_{k+1}\) is exactly \(B_k(h,a)\). Values at next states of zero
conditional probability are unnecessary; any finite completion there
gives the same expectation. The analogous terminal expectation is
\(B_K(h,a)=\mathbb E[\varphi(Y)\mid H_K=h,A_K=a]\).
The two preceding identification theorems supply the causal interpretation
under their respective hypotheses. The algebra itself makes no claim
that a conditional mean of the factual response, without this recursion,
already equals the value of a multi-stage policy.
\end{proof}
\end{annexitem}

\begin{annexitem}{annex:p16.sequential-ipw}{thm:p16-sequential-ipw}{Cancellation on observed paths and the zero-mass complement}
\begin{proof}
If \(p(o)>0\), every prefix event containing that complete assignment
has positive probability. The finite chain rule thus has no null
denominator along that path and gives
\[
 p(o)=b(l_0)\prod_{k=0}^K e_k(a_k\mid h_k)
               \prod_{k<K}r_k(l_{k+1}\mid h_k,a_k)r_K(y\mid h_K,a_K).
\]
In particular every denominator in \(W^\pi(o)\) is strictly positive.
Multiplication by its definition cancels exactly one copy of each of
these positive \(e_k\) factors. The result is (GF). If any policy
factor is zero, this multiplication yields zero on both sides without
an undefined ratio. On an assignment with \(p(o)=0\), the support lemma
gives \(q_\pi(o)=0\); the declared value \(W^\pi(o)=0\) therefore
also proves the mass equality there. It is not an assertion that a
structural intervention at that assignment fails to exist.

Sum the mass equality over all assignments. Normalization of \(q_\pi\)
gives \(\mathbb E_pW^\pi=1\). Multiply the equality by \(\varphi(y)\)
and sum to obtain the weighting identity. The causal conclusions follow
from the appropriate preceding identification theorem. For deterministic
\(g\), the product of its indicator numerators is one exactly when all
recorded actions follow that rule, and zero otherwise.

For a fixed static sequence \(\bar a^*\), the weight is zero outside
that sequence. Multiplication by \(c(\bar a^*)>0\) multiplies both the
weighted-response expectation and the weight expectation by the same
constant. Their ratio stays equal to the original target. A numerator
\(\prod_k s_k(a_k\mid\bar a_{k-1})\), where the \(s_k\) are prescribed
treatment-only probabilities, is a constant at this fixed sequence,
provided that constant is positive. Under a dynamic strategy different
covariate paths can prescribe different sequences and hence different
values of this numerator. Such a factor no longer cancels out of the
sum. The explicit final example below demonstrates failure of the
resulting normalized ratio even under randomized treatment.

Finally, if the policy kernels equal the observed kernels on positive
histories, the observed-law formulation of (RP) is immediate: a policy
action of positive mass has that identical positive observed mass.
On every positive observed assignment each ratio in \(W^\pi\) is one.
The proved mass equality then gives \(q_\pi=p\) there; both masses are
zero elsewhere. This proves preservation of the full population law.
For an explicit failure of unitwise preservation, let a single factual
treatment be fair and its response equal that treatment. A fresh independent
fair policy draw has the same trajectory law, but differs from the factual
treatment, and hence from its factual response, with probability \(1/2\).
\end{proof}
\end{annexitem}

\begin{annexitem}{annex:p16.sequential-source-feedback}{thm:p16-sequential-source-feedback}{Realizing and calculating the source's stipulated table}
\begin{proof}
Take independent \(A_0,U\), both fair bits, \(R_L\) uniform on
\(\{1,\ldots,10\}\), and \(R_A\) uniform on \(\{1,\ldots,5\}\).
Let \(L_0\) be constant and specify the threshold table
\[
\begin{array}{c|cc}
 t(a,u)&u=0&u=1\\\hline
 a=0&1&4\\a=1&3&7
\end{array}
\qquad
 L_1=\mathbf1_{\{R_L>t(A_0,U)\}},\quad
 A_1=\mathbf1_{\{R_A\leq2+2L_1\}},\quad Y=20+80U.
\]
Thus \(\mathbb P(L_1=0\mid a,u)=t(a,u)/10\). These are total
finite structural equations in an acyclic order. Treatment uses its
private randomness conditional on the observed past; \(U,R_L\) are
background variables allowed in the structural bridge. Hence (SE--J)
and consistency hold. Neither treatment occurs in the response equation,
and interventions preserve \(U\), so \(Y^\pi=20+80U\) for every policy
and its mean is \(60\).

Marginalizing the threshold table gives
\[
\begin{array}{c|cc|cc}
 a& p(L_1=0\mid a)&p(L_1=1\mid a)&
       p(U=1\mid a,L_1=0)&p(U=1\mid a,L_1=1)\\\hline
 0&1/4&3/4&4/5&2/5\\
 1&1/2&1/2&7/10&3/10
\end{array}
\]
For example the first conditional posterior is
\((\frac12\frac4{10})/(\frac14)=4/5\); the other three follow by the
same division. Conditional on \(A_0,L_1\), the private \(R_A\) leaves
these posteriors unchanged on either \(A_1\) stratum. Multiplying them
by \(80\) and adding \(20\) gives the four means \(84,52,76,44\).
Also \(p(A_1=1\mid L_1=0)=2/5\) and
\(p(A_1=1\mid L_1=1)=4/5\). Multiplying \(p(A_0)=1/2\), each
covariate probability and each action probability gives respectively
the eight counts in the body divided by \(32000\).

This specifies the full joint response distribution, not only its means.
Multiplying each row mass by the indicated posterior and its complement
gives the complete table, with every entry divided by \(32000\):
\[
\begin{array}{ccc|rr}
 A_0&L_1&A_1&Y=20&Y=100\\\hline
 0&0&0&480&1920\\0&0&1&320&1280\\
 0&1&0&1440&960\\0&1&1&5760&3840\\
 1&0&0&1440&3360\\1&0&1&960&2240\\
 1&1&0&1120&480\\1&1&1&4480&1920
\end{array}
\]
Each response column sums to \(16000\). The cells are positive,
their row sums and response means are the body table's entries, and
their total mass is one. Every treatment option at every observed
history is positive.

For never treating, the g-formula sums \(84(1/4)+52(3/4)=60\).
For always treating, it sums \(76(1/2)+44(1/2)=60\). More generally
\(A_1\) does not enter any of the conditional means, so summing any
second-stage policy gives the same conditional contribution before that
action. At either first-stage choice, the remaining covariate sum is
\(60\); any first-stage randomization also gives \(60\).

The nonstabilized static inverse weights are
\(1/[\frac12 p(A_1=a_1\mid L_1)]\). On the two never-treating rows
they are \(10/3\) and \(10\). Multiplying their counts \(2400,2400\)
gives \(8000,24000\), whose weighted response mean is
\((8000\cdot84+24000\cdot52)/32000=60\).
On the always-treating rows the weights are \(5\) and \(5/2\);
counts \(3200,6400\) become \(16000,16000\), giving mean \(60\).
The total weighted count is \(32000\) in either static sequence,
so the normalization agrees with \(\mathbb E W^g=1\).

Without those weights, the always-treating mean is
\((3200\cdot76+6400\cdot44)/9600=164/3\), whereas the never-treating
mean is \((2400\cdot84+2400\cdot52)/4800=68\). Their difference
is \(-40/3\). Within \(L_1=0\), the corresponding difference is
\(76-84=-8\), and within \(L_1=1\) it is \(44-52=-8\).
Here conditioning on \(L_1\), affected by both \(A_0\) and \(U\),
changes the distribution of the response's background cause \(U\)
between first-stage treatment groups, as the posterior table explicitly
shows. The source's treatment-confounder feedback phenomenon is therefore
realized by a complete finite null model. Its hypothetical counts are
used as exact population proportions, without an empirical or sampling
interpretation.
\end{proof}
\end{annexitem}

\begin{annexitem}{annex:p16.sequential-feedback-effects}{thm:p16-sequential-feedback-effects}{Three policy calculations with nonzero effects}
\begin{proof}
The three original independent bits give \(p(A_0=a)=1/2\),
\[
 p(L_1=1\mid A_0=a)=1/4+a/2,
 \qquad p(A_1=1\mid L_1=l)=1/4+l/2.
\]
Their products give the body table: for \(A_0=0\), the four masses
are \(9,3,1,3\) divided by \(32\); for \(A_0=1\) they are
\(3,1,3,9\) divided by \(32\). The response is the indicated
deterministic value \(l+2a_1\). These masses sum to one, and every
history-action combination has positive probability. Treatment's private
bit mechanism meets the structural bridge, so all causal premises hold.

After static replacement, \(L_1=a_0\oplus E_L\) and
\(Y=l+2a_1\). Consequently the four means in sequence order
\((0,0),(0,1),(1,0),(1,1)\) are
\(1/4,9/4,3/4,11/4\). This is also the g-formula sum
\(\sum_l(l+2a_1)p(l\mid a_0)\).
For the deterministic second-stage rule \(a_1=l\), the terminal
backward value is \(V_1(a,l)=3l\). Hence
\(B_0(a)=3(1/4+a/2)=3/4+3a/2\).
For the stochastic second-stage rule in the body,
\[
 V_1(a,l)=l+2(1/5+3l/5)=2/5+11l/5.
\]
Integrating this with \(p(l\mid a)\) gives
\(2/5+(11/5)(1/4+a/2)=19/20+11a/10\).
This computes all the g-formula and backward-recursion values directly.

For completeness, take any one of these second-stage policies \(\pi_1\)
and fix \(A_0=a\). On row \((a,l,b)\), the weight is
\(2\pi_1(b\mid a,l)/e_1(b\mid l)\). Multiplying its positive
row mass \(\frac12p(l\mid a)e_1(b\mid l)\) gives exactly
\(p(l\mid a)\pi_1(b\mid a,l)\). Summing the resulting response
contributions \((l+2b)p(l\mid a)\pi_1(b\mid a,l)\) gives the same
three calculations just made. Rows with another \(A_0\) have weight
zero, and rows with an excluded deterministic second action also have
weight zero. Thus the weighting is fully specified in every observed row.

In the unweighted table the never-treating mean is \(1/10\), since its
two row counts are \(9,1\) with responses \(0,1\). The always-treating
mean is \((1\cdot2+9\cdot3)/10=29/10\). Their difference is
\(14/5\), whereas the policy difference is \(11/4-1/4=5/2\).
Within either fixed \(l\), the response contrast is
\((l+2)-l=2\). Fixing the covariate in that comparison removes the
additional \(1/2\) response change caused through \(A_0\to L_1\to Y\).
The same \(L_1\) also determines the later treatment probability.
The graph therefore has feedback in time \(A_0\to L_1\to A_1\),
without a directed cycle, and the calculations exhibit its consequences.
\end{proof}
\end{annexitem}

\begin{annexitem}{annex:p16.sequential-failures}{thm:p16-sequential-failures}{Explicit equal-law and exchangeability counterexamples}
\begin{proof}
In the first pair, take the two equally likely values of a treatment bit,
constant covariates, and the structural equations stated in the body.
Both models produce exactly \((A_0,A_1,Y)=(0,0,0),(1,0,0)\), each with
mass \(1/2\). Every fixed-world response and covariate is constant, so
all the (SE--J) independences hold. The first action zero is supported.
At its reached second-stage history, the policy asks for \(A_1=1\),
whose observed probability is zero. The structural replacement still
exists: it yields zero in the first model and one in the second. Any
assigned observed conditional mean at this absent action is merely a
version choice and cannot distinguish the models. For the policy \((0,0)\)
both stages are supported and the identified mean is zero.

For the second pair, conditional on observed \(A=U=a\), the first
response equation reduces to \(Y=E\), exactly the second model's
equation. Thus in both models \(p(a,0)=\frac12\frac34=3/8\) and
\(p(a,1)=\frac12\frac14=1/8\). Every one of these four cells is
positive, and either action has probability \(1/2\). Under \(A=a\)
replacement, the first response becomes \(a\oplus U\oplus E\).
Since \(U\) is an unchanged fair bit independent of \(E\), this is
fair. The second remains \(E\), with success probability \(1/4\).
In the first model, conditional on factual \(A=a\), the success
probability of \(Y[a]\) is \(1/4\), while conditional on factual
\(A=1-a\) it is \(3/4\). Thus (SE--Y) fails explicitly, although
the support condition holds. The observed g-formula and weighting both
return \(1/4\), displaying the missing causal premise.

For the third example take independent fair bits \(U,V,E\), constant
\(L_0\), and the total equations
\[
 A_0=U\oplus V,\qquad L_1=U,\qquad A_1=E,
 \qquad Y=V\oplus A_1.
\]
The order is acyclic with background \((U,V)\) allowed to affect both
\(A_0\) and later variables. Every observed \((A_0,L_1,A_1)\) triple
has mass \(1/8\); given such a triple, \(Y=A_0\oplus L_1\oplus A_1\).
Thus \(p(A_0=a)=1/2\) and
\(p(A_1=b\mid A_0=a,L_1=l)=1/2\), proving every required treatment
positivity condition.

For every static \((a_0,a_1)\),
\(L_1[a_0]=U\) and \(Y[a_0,a_1]=V\oplus a_1\). At the first stage
each is separately independent of \(A_0=U\oplus V\): conditioning on
either value of \(A_0\) leaves each of \(U,V\) fair. Their pair,
however, determines \(A_0\) after the known bit \(a_1\) is accounted
for. In particular for \(a_1=0\),
\[
 \mathbb P(U=0,V=0\mid A_0=0)=1/2,
 \qquad\mathbb P(U=0,V=0\mid A_0=1)=0.
\]
So first-stage joint exchangeability fails. At the second stage,
\(Y[a_0,a_1]\) is a function of \(U,V\) and is independent of the
private fair \(E\), also conditional on \(H_1\). Thus the static
outcome-specific exchangeability conditions hold at both times for all
four sequences. Each static outcome is fair, and the static g-formula
also gives \(1/2\).

Now prescribe \(A_0=0\) followed by \(A_1=L_1\). Its true response
is \(Y^g=V\oplus U\), a fair bit equal to the factual \(A_0\).
Consequently its first-stage exchangeability fails. In the observed
stratum \(A_0=0\), one has \(V=U=L_1\); on adherence at the second
stage the observed response is therefore zero. The observed g-formula
gives \(0\cdot\frac12+0\cdot\frac12=0\), whereas
\(\mathbb E[Y^g]=1/2\). This proves the strict distinctions claimed
between static outcome conditions, separate marginal conditions and
the joint premise used for dynamic identification. All counterfactuals
in these examples come from actual equation replacement, so consistency
holds pointwise throughout.
\end{proof}
\end{annexitem}

\begin{annexitem}{annex:p16.sequential-weight-boundaries}{thm:p16-sequential-weight-boundaries}{The exact normalized targets and a fitted-denominator failure}
\begin{proof}
In the example, \(p(L=0,A=0)=1/8\),
\(p(L=0,A=1)=3/8\), \(p(L=1,A=0)=1/8\), and
\(p(L=1,A=1)=3/8\). The equation \(Y=L\) is unaffected by treatment,
so its mean under \(A=L\) is \(1/2\). Treatment is independent of
the covariate and all potential outcomes, and both options have positive
conditional probability. The causal and support assumptions hold.

The treatment-only numerator equals the actual conditional denominator,
so the proposed stabilized weight \(p(A)/p(A\mid L)\) is one in every
row. Restricting to adherers leaves rows \((0,0)\) and \((1,1)\)
with masses \(1/8,3/8\) and responses zero and one. Their normalized
mean is \((3/8)/(1/8+3/8)=3/4\). The correct policy weights are
\(4\) in the first adherent row and \(4/3\) in the second, and zero
elsewhere. The resulting masses are \(1/2,1/2\), giving total weight
one and weighted mean \(1/2\). This computes the target change due
to dynamic stabilization without an appeal to an estimation argument.

For identically distributed observations, finite linearity gives
\[
 \mathbb E\left[\frac1n\sum_{i=1}^nW^\pi(O_i)Y_i\right]
 =\frac1n\sum_{i=1}^n\mathbb E_p[W^\pi(O)Y]
 =\psi(\pi).
\]
This uses the true weights and the already proved identification identity;
no independence across observations enters this expectation calculation.
It is an assertion about the unnormalized sample average, not the sample
ratio dividing by an estimated or random mean weight.

To demonstrate failure with incorrect denominators even before fitting,
take one randomized fair treatment \(A\), no covariates and \(Y=A\).
The target under \(A=1\) is one. Replace its true denominator \(1/2\)
by the fixed proposed value \(3/4\), with complementary proposed
probability \(1/4\) for zero. The proposed weighted-response expectation
is \(\mathbb E[\mathbf1_{\{A=1\}}Y/(3/4)]=(1/2)/(3/4)=2/3\).
Thus even a normalized positive proposed treatment kernel need not give
the identity. A fitted random kernel requires additional analysis of
its relation to the actual assignment probabilities and the sample.
Population algebra alone proves neither consistency nor an asymptotic
distribution for such fitted procedures.
\end{proof}
\end{annexitem}
